\chapter{Cómo usar este manual}

Este manual explica, de principio a fin, cómo se construye, se pone en marcha, se usa y se mantiene un poste de iluminación \dato{producto}. Está escrito para el técnico que instala postes y que, si hace falta, también modifica el código. No supone que conozcas WLED, React, KiCad ni LaTeX: cada término aparece en el glosario de este capítulo.

\section{Qué es un poste LUMA}

Un poste LUMA tiene dos caras de luces LED direccionables (\dato{tipoLed}), con \dato{ledsPorCara} píxeles por cara (\dato{ledsTotal} en total). Las dos caras muestran siempre lo mismo, píxel por píxel, en todos los efectos. Una placa propia, la \dato{pcbVersion}, reparte la energía de una fuente de 12~V y lleva una \dato{placaEsp} que ejecuta el firmware libre WLED. El poste crea su propia red Wi-Fi, «\dato{redPropiaSsid}», y sirve la aplicación LUMA en \texttt{http://\dato{redPropiaIp}}: cualquier teléfono que se una a esa red controla las luces sin internet ni computador.

\begin{figure}[H]\centering
\resizebox{\linewidth}{!}{% Diagrama de bloques del sistema completo
\begin{tikzpicture}[
  font=\sffamily\small, node distance=8mm and 10mm,
  bloque/.style={draw=grafito, fill=marfil, rounded corners=2pt, align=center, minimum height=10mm, inner sep=5pt},
  placa/.style={draw=acento, dashed, rounded corners=4pt, inner sep=6pt},
  cara/.style={draw=grafito, fill=lima!45, rounded corners=2pt, align=center, minimum height=10mm, minimum width=24mm},
  energia/.style={-{Stealth[length=2.2mm]}, very thick, rojoaviso!80!black},
  datos/.style={-{Stealth[length=2.2mm]}, thick, azulnota},
  red/.style={{Stealth[length=2.2mm]}-{Stealth[length=2.2mm]}, thick, dashed, acento},
]
  \node[bloque] (red120) {Red eléctrica\\120 V AC};
  \node[bloque, right=of red120] (fuente) {Fuente 12 V\\{\scriptsize\dato{fuente}}};

  % Dentro de la PCB
  \node[bloque, below=14mm of fuente, xshift=-28mm] (lm) {LM2596\\salida \dato{reguladorSalida}};
  \node[bloque, below=of lm] (esp) {ESP32-S3 N16R8\\WLED + LUMA};
  \node[bloque, right=22mm of lm] (fus) {Fusibles\\F6 (cara A) · F7 (cara B)};
  \node[bloque, below=of fus] (bornes) {Borneras A1 · B1\\(+12 V, datos, GND)};
  \node[placa, fit=(lm)(esp)(fus)(bornes), label={[font=\sffamily\small\bfseries\color{acento}, anchor=south east]north east:PCB\_V2}] (pcb) {};

  \node[cara, below=14mm of bornes, xshift=-16mm] (caraA) {Cara A\\\dato{ledsPorCara} LED};
  \node[cara, below=14mm of bornes, xshift=16mm] (caraB) {Cara B\\\dato{ledsPorCara} LED};

  \node[bloque, left=36mm of esp] (tel) {Teléfono\\con LUMA abierta};

  \draw[energia] (red120) -- (fuente);
  \draw[energia] (fuente.south) -- ++(0,-5mm) -| (lm.north);
  \draw[energia] (fuente.south) -- ++(0,-5mm) -| (fus.north);
  \draw[energia] (lm) -- node[right, font=\scriptsize\sffamily]{5 V} (esp);
  \draw[energia] (fus) -- (bornes);
  \draw[datos] (esp.east) -- node[below, font=\scriptsize\sffamily, align=center]{datos\\A y B} (bornes.west);
  \draw[energia] (bornes.south) -- ++(0,-5mm) -| (caraA.north);
  \draw[energia] (bornes.south) -- ++(0,-5mm) -| (caraB.north);
  \draw[red] (tel) -- node[above, font=\scriptsize\sffamily, align=center]{Wi-Fi «\dato{redPropiaSsid}»\\http://\dato{redPropiaIp}} (esp);

  \matrix[below=6mm of caraA.south west, anchor=north west, column sep=4pt, font=\scriptsize\sffamily, xshift=-30mm] {
    \draw[energia] (0,0) -- (6mm,0); & \node{energía}; &
    \draw[datos] (0,0) -- (6mm,0); & \node{datos LED}; &
    \draw[red] (0,0) -- (6mm,0); & \node{Wi-Fi}; \\
  };
\end{tikzpicture}
}
\caption{Visión general del sistema.}\label{fig:bloques}
\end{figure}

\section{Cómo está organizado}

\begin{description}[style=nextline, leftmargin=0pt]
\item[Parte I · Hardware] Materiales, la placa, el esquema eléctrico, el ensamble, el cableado del poste y el consumo.
\item[Parte II · Firmware] Preparar el computador, grabar WLED en una ESP32-S3 nueva y configurarlo para el poste.
\item[Parte III · App LUMA] Instalar la aplicación dentro del poste, usarla y probar el poste antes de entregarlo.
\item[Parte IV · El código] Cómo está hecho LUMA, cómo hacer los cambios más comunes y cómo probarlos.
\item[Parte V · Operación] Solución de problemas, actualizaciones, respaldos y limitaciones conocidas.
\item[Anexos] Tabla de pines, configuración comentada, efectos, hoja de verificación, referencias y licencias.
\end{description}

Si vas a montar un poste nuevo, sigue las partes I, II y III en orden y termina con la hoja de verificación (anexo~D). Si vas a cambiar el código, lee primero el capítulo de arquitectura (parte IV).

\section{Convenciones}

Los recuadros señalan lo que no debe pasarse por alto:

\advertencia{Un paso que puede dañar el equipo o causar un accidente si se hace mal. Léelo antes de seguir.}

\nota{Información útil: una aclaración, un atajo o el porqué de un paso.}

\pendiente{fusibles}{Un dato que el equipo todavía debe confirmar. Mientras exista este recuadro, el manual no da ese dato por cierto.}

Los comandos que se escriben en una terminal aparecen así:

\begin{comando}
npm run build:poste
\end{comando}

Los nombres de botones y pantallas de LUMA se escriben entre comillas angulares, por ejemplo «Mi montaje». Las rutas de archivos y los nombres técnicos van en \texttt{letra de máquina}. Donde falta una foto real del montaje aparece un marco «FOTO PENDIENTE» con la descripción de lo que hay que fotografiar.

\section{Glosario}

\begin{description}[leftmargin=0pt, itemsep=0.3em]
\item[AP (punto de acceso)] Red Wi-Fi que crea el propio poste («\dato{redPropiaSsid}»). Funciona sin router ni internet.
\item[cfg.json] Archivo de configuración de WLED guardado en la memoria de la ESP32: pines, cantidad de LED, red, horarios.
\item[Escena] Combinación guardada de colores, efecto, segmentos y brillo. LUMA la guarda en el teléfono y, si se marca, en el poste (como preset de WLED).
\item[ESP32-S3] Microcontrolador con Wi-Fi que controla las luces. La placa usada es una \dato{placaEsp}.
\item[Firmware] Programa grabado en la ESP32. Aquí es WLED versión \dato{wledVersion}.
\item[GPIO] Pin de entrada/salida de la ESP32. Cada cara de LED recibe sus datos por un GPIO.
\item[I2S y RMT] Dos periféricos internos de la ESP32 que WLED usa para generar la señal de los LED. En este poste la cara A usa RMT y la cara B debe usar I2S (ver capítulo~\ref{cap:wled}).
\item[Salida de LED] Cada pin de datos configurado en WLED, con su LED inicial («Start») y su cantidad. En el poste, las dos salidas empiezan en el LED 0 para que la cara B repita a la cara A.
\item[Inyección] Alimentar la tira también por el extremo lejano para que no caiga la tensión a lo largo del poste.
\item[OTA] Actualización del firmware por Wi-Fi, sin cable USB.
\item[Píxel] Cada punto de luz controlable. En las tiras WS2811 de 12~V, un chip controla un píxel.
\item[Preset] Escena guardada dentro de WLED; LUMA usa presets para las escenas del poste, el modo automático y el horario.
\item[Segmento] Tramo de LED con su propio efecto y colores. LUMA admite hasta 16.
\item[WLED] Firmware libre para controlar LED direccionables desde una ESP32 (\url{https://kno.wled.ge}).
\item[WS2811] Tipo de chip de los píxeles. Recibe datos en serie: el primer píxel toma su color y pasa el resto al siguiente.
\end{description}
